\begin{abstract}
\noindent Does the Hurst exponent $H$ of a cross-country flight depend on the glider? We compare
\UniFunnelParaOpen{} paraglider flights and \UniFunnelHangOpen{} hang glider flights of the French archive, all
declared open with an opening above \UniMinOpening, on the mean squared displacement (MSD) of the whole
continuous stretches of the track. On a fixed population of flights, $H$ is \UniHFirstAllAfixed{} for paragliders and
\UniHSecondAllAfixed{} for hang gliders, a difference of \UniDHAllAfixed{} with a 90\% band \UniDHBandAllAfixed. The
difference is positive in \UniForestPos{} of \UniForestN{} populations and designs, between \UniForestMin{} and
\UniForestMax, in every terrain class (the class of a flight is the class of the median ground elevation under its
points) and in the Alps. $H$ also falls from plains to mountains in both disciplines, by \UniHFirstPlainsAfixed{} to \UniHFirstHighAfixed{} for
paragliders and \UniHSecondPlainsAfixed{} to \UniHSecondHighAfixed{} for hang gliders. The gap between the gliders is
of the order of the gap between two terrains: $H$ is universal to a few hundredths, and not to the precision of the
bands.
\end{abstract}
